\hypertarget{gpu-acceleration-with-cuda-and-python}{%
\chapter{GPU Acceleration with CUDA and
Python}\label{gpu-acceleration-with-cuda-and-python}}

When the CPU's 8-16 cores aren't enough, we turn to the GPU, which can
have thousands of cores.

\hypertarget{the-cuda-architecture}{%
\subsection{72.1 The CUDA Architecture}\label{the-cuda-architecture}}

CUDA (Compute Unified Device Architecture) is NVIDIA's platform for
parallel computing. * \textbf{Kernels}: Small functions that run on the
GPU. * \textbf{Memory Transfer}: Data must be moved from Host (CPU RAM)
to Device (GPU RAM) before processing.

\hypertarget{interfacing-with-python-cupy-and-numba}{%
\subsection{72.2 Interfacing with Python: CuPy and
Numba}\label{interfacing-with-python-cupy-and-numba}}

\begin{itemize}
\tightlist
\item
  \textbf{CuPy}: A NumPy-compatible library that runs on the GPU.
\item
  \textbf{Numba \passthrough{\lstinline!@cuda.jit!}}: A JIT compiler
  that translates Python functions directly into PTX (GPU machine code).
\end{itemize}

\begin{center}\rule{0.5\linewidth}{0.5pt}\end{center}

To truly master Python, one must understand how it compares to its peers
and where it is headed in the next decade.

\hypertarget{chapter-75-comparative-analysis-python-vs.-c-vs.-rust}{%
\chapter{Chapter 75: Comparative Analysis: Python vs.~C++
vs.~Rust}\label{chapter-75-comparative-analysis-python-vs.-c-vs.-rust}}

Choosing the right tool for the job requires an objective look at the
trade-offs between these three dominant languages.

\hypertarget{performance-vs.-productivity}{%
\subsection{75.1 Performance
vs.~Productivity}\label{performance-vs.-productivity}}

\begin{longtable}[]{@{}llll@{}}
\toprule
Feature & Python & C++ & Rust \\
\midrule
\endhead
\textbf{Execution Speed} & Moderate (Interpreter) & Extreme (AOT) &
Extreme (AOT) \\
\textbf{Development Speed} & High & Low & Moderate \\
\textbf{Memory Safety} & Managed (GC) & Manual (Unsafe) & Managed
(Borrow Check) \\
\textbf{Concurrency} & Cooperative/Preemptive & Preemptive &
Preemptive \\
\bottomrule
\end{longtable}

\hypertarget{the-godhood-perspective}{%
\subsection{75.2 The ``Godhood''
Perspective}\label{the-godhood-perspective}}

\begin{itemize}
\tightlist
\item
  \textbf{Python}: Best for high-level orchestration, rapid prototyping,
  and data science where developer time is more expensive than CPU time.
\item
  \textbf{C++}: Best for legacy systems, game engines, and scenarios
  requiring absolute control over hardware.
\item
  \textbf{Rust}: The modern choice for systems programming, providing
  C++ performance with guaranteed memory safety.
\end{itemize}

\begin{center}\rule{0.5\linewidth}{0.5pt}\end{center}
